\section{Dependence and the score subspace}
For a fixed prompt, let $p_\theta(y)>0$ be a finite-support differentiable policy,
$s(y)=\nabla_\theta\log p_\theta(y)$, and $r(y)$ a parameter-independent reward.
Let an exchangeable group $(Y_1,\ldots,Y_K)\sim Q_\theta$ have marginal $p_\theta$.
The intended gradient is $g=\E_p[sr]$. At the generating checkpoint, the sequence
RLOO estimator is
\begin{equation}
\widehat g_Q=\frac1K\sum_i\left(r_i-\frac1{K-1}\sum_{j\ne i}r_j\right)s_i.
\label{eq:rloo}
\end{equation}
Advantages are detached, and sequence scores sum all response-token log-probability
gradients, including emitted EOS. A capped response remains an outcome. This
estimand excludes clipping, off-policy epochs, length normalization and KL rewards,
axes along which practical implementations differ
\citep{grpo2026ustatistic,trl}.

Define $(Tf)(y)=\E[f(Y_2)\mid Y_1=y]$ and $L=I-T$. Exchangeability gives the known
pairwise identity
\begin{equation}
g_Q=\E_p[sLr]=\tfrac12\E_Q[(s_1-s_2)(r_1-r_2)].\label{eq:operator}
\end{equation}
In $L_2(p)$, $T$ is self-adjoint and contractive, and $T1=1$.
Its Dirichlet form is $\inner{f}{Lf}=\E[(f(Y_1)-f(Y_2))^2]/2\ge0$.
Extendibility to $K$ exchangeable members gives
$-I/(K-1)\preceq T\preceq I$ on centered functions. A reversible pair matrix
alone does not guarantee extendibility for an arbitrary larger $K$.

Let $(Sa)(y)=s(y)^Ta$, $F=S^*S$, $\V=\operatorname{range}(S)$ and
$\Pi=SF^\dagger S^*$. Decompose centered reward as
$r_0=r-\E r=v+e$ with $v=\Pi r_0$ and $e\perp\V$. Then
\begin{equation}
D_Q(r):=g^TF^\dagger g_Q
=\inner{v}{Lv}-\inner{v}{Te}.\label{eq:decomposition}
\end{equation}
The first term is nonnegative; the second describes residual interference.

\begin{theorem}[Universal local compatibility]
For a fixed positive finite-support policy and exchangeable sampler, the following
are equivalent: (i) $D_Q(r)\ge0$ for every parameter-independent reward;
(ii) $T\V\subseteq\V$; (iii) $\Pi T(I-\Pi)=0$;
(iv) one reward-independent positive self-adjoint operator $A$ on $\V$ satisfies
$\Pi Lr=A\Pi r$ for every reward. In that case $A=\Pi L\Pi|_\V$.
\end{theorem}
\begin{proof}
In $\V\oplus\V^\perp$, write $L=\left(\begin{smallmatrix}A&B\\B^*&C\end{smallmatrix}\right)$.
The slope is $\inner{v}{Av}+\inner{v}{Be}$. If $B=0$, positivity of $L$ gives
nonnegative slope. Otherwise choose a nonzero cross term and scale and sign $e$
to make the slope negative. On finite support, an affine reward transformation
maps this example into $[0,1]$ without changing its sign. Self-adjointness makes
the off-diagonal blocks vanish together, yielding the equivalent conditions.
\end{proof}

This is a natural-gradient statement, not an AdamW ascent theorem. A saturated
score subspace has $e=0$ and cannot exhibit natural reversal. One observed acute
gradient pair cannot establish reward-independent preconditioning. Shared model
parameters also mean positive prompt-specific scaling may rotate an aggregate update,
as finite-group objective analyses of prompt weighting also emphasize
\citep{softmaxgrpo2026}.

\subsection{Stratified reweighting and lattice filtering}
Fix a decoder $Y=G_\theta(U)$ with uniform $U\in[0,1)$ and a declared token ordering.
Let $\mathcal C f=f\circ G_\theta$. For the projection $P_B$ onto means in $K$
equal latent bins, define $H_B=\mathcal C^*P_B\mathcal C$.
One independent sample per bin, followed by random labeling, gives on centered functions
\begin{equation}
L_{\rm strat}=I+\frac{H_B}{K-1},\qquad 0\preceq H_B\preceq I.\label{eq:strat}
\end{equation}
Thus if $\eta=\|v\|_p^2/\|r_0\|_p^2>0$,
\begin{equation}
\frac{D_{\rm strat}(r)}{\|v\|_p^2}
\ge\frac{2K-1-\eta^{-1/2}}{2(K-1)}.\label{eq:bound}
\end{equation}
In particular, $\eta>1/(2K-1)^2$ suffices for strictly positive natural slope;
falling below the threshold only permits reversal.

For lattice sampling, define the orbit projection
$P_Of(u)=K^{-1}\sum_{j=0}^{K-1}f((u+j/K)\bmod1)$ and
$H_O=\mathcal C^*P_O\mathcal C$. Then
\begin{equation}
L_{\rm lattice}=\frac K{K-1}(I-H_O).\label{eq:lattice}
\end{equation}
An orbit-invariant lifted reward receives no RLOO signal. Compression matters:
$H_B,H_O$ are positive contractions and need not themselves be projections on sequence space.
